% 10_complete_paper.tex
% Combined Digital Humanities + Siggraph Paper Template
% Overleaf compatible: YES (article class with optional acmart for Siggraph)

% ============================================================
% DUAL-FORMAT TEMPLATE
% ============================================================
% This file can be compiled in two modes:
% 1. DH Mode (default): Single-column article with critical apparatus
% 2. Siggraph Mode: Two-column acmart format
%
% Switch by changing \dhformattrue / \dhformatfalse below
% ============================================================

% Format switch - set to true for DH, false for Siggraph
\newif\ifdhformat
\dhformattrue  % <-- CHANGE THIS TO SWITCH FORMATS

\ifdhformat
% ============================================================
% DIGITAL HUMANITIES FORMAT (Single-column article)
% ============================================================
\documentclass[12pt,a4paper]{article}
\usepackage[utf8]{inputenc}
\usepackage[T1]{fontenc}
\usepackage{lmodern}
\usepackage[english,french,latin,greek]{babel}
\usepackage{hyperref}
\usepackage{cleveref}
\usepackage{geometry}
\usepackage{fontspec}
\usepackage{polyglossia}
\usepackage{reledmac}
\usepackage{reledpar}
\usepackage{parallel}
\usepackage{verse}
\usepackage{graphicx}
\usepackage{subcaption}
\usepackage{booktabs}
\usepackage{multirow}
\usepackage{makecell}
\usepackage{longtable}
\usepackage{float}
\usepackage{caption}
\usepackage{siunitx}
\usepackage{tikz}
\usepackage{pgfplots}
\usepackage{amsmath}
\usepackage{amssymb}
\usepackage{mathtools}
\usepackage{physics}
\usepackage{algorithm2e}
\usepackage{algpseudocode}
\usepackage{listings}
\usepackage{xcolor}
\usepackage{footmisc}
\usepackage{marginnote}

\setmainlanguage{english}
\setotherlanguages{french,latin,greek}

% Critical apparatus setup
\Xarrangement[A]{paragraph}
\Xarrangement[B]{paragraph}
\linenumincrement*{5}
\firstlinenum*{1}

\geometry{marginparwidth=3.5cm, marginparsep=0.5cm}
\usepackage[stable]{footmisc}
\setlength{\footnotemargin}{1.5em}

\else
% ============================================================
% SIGGRAPH FORMAT (Two-column acmart)
% ============================================================
\documentclass[sigconf,review,anonymous]{acmart}
\usepackage[utf8]{inputenc}
\usepackage[T1]{fontenc}
\usepackage{lmodern}
\usepackage[english]{babel}
\usepackage{hyperref}
\usepackage{cleveref}
\usepackage{graphicx}
\usepackage{subcaption}
\usepackage{booktabs}
\usepackage{amsmath}
\usepackage{amssymb}
\usepackage{mathtools}
\usepackage{physics}
\usepackage{algorithm2e}
\usepackage{algpseudocode}
\usepackage{listings}
\usepackage{xcolor}

\acmConference[SIGGRAPH '26]{ACM SIGGRAPH 2026 Conference Papers}{August 2026}{Denver, CO, USA}
\acmYear{2026}
\acmISBN{978-1-4503-XXXX-X/26/08}
\acmDOI{10.1145/XXXXXXX.XXXXXXX}
\linenumbers

\fi

% ============================================================
% COMMON SETUP (both formats)
% ============================================================
\hypersetup{
    colorlinks=true,
    linkcolor=blue,
    citecolor=red,
    urlcolor=cyan
}

\SetAlgoLined
\SetKwInOut{Input}{Input}
\SetKwInOut{Output}{Output}

\lstset{
    basicstyle=\ttfamily\small,
    keywordstyle=\color{blue},
    commentstyle=\color{green!60!black},
    stringstyle=\color{red},
    numbers=left,
    numberstyle=\tiny,
    stepnumber=1,
    numbersep=5pt,
    frame=single,
    breaklines=true,
    showstringspaces=false
}

\pgfplotsset{compat=1.18}

\usetikzlibrary{
    arrows.meta, positioning, shapes.geometric, shapes.misc,
    calc, decorations.pathreplacing, decorations.markings,
    matrix, fit, backgrounds, shadows, spy, external
}

% Theorem environments
\newtheorem{theorem}{Theorem}[section]
\newtheorem{lemma}[theorem]{Lemma}
\newtheorem{proposition}[theorem]{Proposition}
\newtheorem{corollary}[theorem]{Corollary}
\newtheorem{definition}[theorem]{Definition}
\newtheorem{example}[theorem]{Example}
\newtheorem{remark}[theorem]{Remark}

% ============================================================
% BIBLIOGRAPHY
% ============================================================
\ifdhformat
\usepackage[
    backend=biber,
    style=authoryear-comp,
    sorting=nyt,
    maxcitenames=2,
    maxbibnames=99,
    doi=true,
    url=true,
    isbn=false
]{biblatex}
\addbibresource{references.bib}
\else
% acmart uses BibTeX with ACM style
\bibliographystyle{ACM-Reference-Format}
\fi

% ============================================================
% PAPER METADATA
% ============================================================
\title{Unified Framework for Neural Rendering and Critical Text Edition}

\ifdhformat
\author{Jane Smith\thanks{Equal contribution} \and Robert Jones\thanks{Equal contribution} \and Min-jun Lee}
\affiliation{%
    \institution{Institute for Digital Humanities}
    \city{Berlin}
    \country{Germany}
}
\affiliation{%
    \institution{Computer Graphics Lab}
    \city{Tokyo}
    \country{Japan}
}
\date{\today}
\else
% Anonymous for review
\author{Anonymous Author(s)}
\affiliation{\institution{Anonymous Institution}}
\fi

% ============================================================
% DOCUMENT START
% ============================================================
\begin{document}

\maketitle

\begin{abstract}
We present a unified framework bridging neural rendering techniques
from computer graphics with critical edition methodologies from
Digital Humanities. Our approach adapts differentiable rendering
pipelines for manuscript image analysis, enabling automated
collation, variant detection, and stemma reconstruction. We
demonstrate applications to both Siggraph-style rendering research
and DH critical editions, with a shared mathematical foundation
in volumetric light transport and textual variation modeling.
\end{abstract}

\ifdhformat
% DH-specific: Table of contents
\tableofcontents
\newpage
\fi

\ifdhformat
\else
\begin{CCSXML}
<ccs2012>
<concept>
<concept_id>10010147.10010371.10010352.10010379</concept_id>
<concept_desc>Computing methodologies~Rendering</concept_desc>
<concept_significance>500</concept_significance>
</concept>
<concept>
<concept_id>10010405.10010432.10010441</concept_id>
<concept_desc>Applied computing~Digital libraries and archives</concept_desc>
<concept_significance>400</concept_significance>
</concept>
</ccs2012>
\end{CCSXML}

\ccsdesc[500]{Computing methodologies~Rendering}
\ccsdesc[400]{Applied computing~Digital libraries and archives}

\keywords{neural rendering, critical edition, differentiable rendering, stemma codicum, manuscript analysis}

\begin{teaserfigure}
\centering
\fbox{\parbox{0.9\linewidth}{\centering [Teaser: Neural collation pipeline showing manuscript $\to$ features $\to$ variants $\to$ stemma]}}
\caption{Our neural collation pipeline: manuscript images are processed through a differentiable renderer to detect textual variants and reconstruct the stemma codicum.}\label{fig:teaser}
\end{teaserfigure}
\fi

% ============================================================
% 1. INTRODUCTION
% ============================================================
\section{Introduction}
\label{sec:intro}

Computer graphics and Digital Humanities share a common challenge:
\emph{reconstructing structure from noisy observations}. In graphics,
we infer 3D scenes from 2D images via differentiable rendering. In DH,
we reconstruct archetypal texts from variant manuscript witnesses
via critical edition. This paper establishes a formal connection.

\ifdhformat
\marginnote{Key insight: Both fields solve inverse problems using differentiable optimization.}
\else
\fi

\subsection{Background: Neural Rendering}
\label{sec:bg-rendering}

Neural rendering \parencite{mildenhall2021nerf,tewari2022advances} represents
scenes as continuous volumetric functions $F_\theta: (\vb{x}, \omega) \mapsto (c, \sigma)$
mapping position and direction to color and density. Training minimizes
photometric loss:

\begin{equation}\label{eq:nerf-loss}
\mathcal{L} = \sum_{r \in \mathcal{R}} \norm{C(r) - \hat{C}(r)}^2
\end{equation}

where $C(r)$ is the rendered color of ray $r$ and $\hat{C}(r)$ the target.

\subsection{Background: Critical Edition}
\label{sec:bg-edition}

A critical edition constructs a constituted text from witnesses
$W = \{w_1, \dots, w_n\}$ by identifying variants at variation units
$V = \{v_1, \dots, v_m\}$. The stemma codicum represents the
genealogical relationships as a directed acyclic graph.

\begin{definition}[Variation Unit]\label{def:varunit}
A variation unit $v$ is a locus where witnesses disagree,
with readings $r_v(w_i) \in \mathcal{R}_v$ for each witness $w_i$.
\end{definition}

\subsection{Our Contribution}
\label{sec:contrib}

We make three contributions:
\begin{enumerate}
    \item A differentiable collation model treating manuscript images
          as ``renderings'' of an underlying archetypal text.
    \item A variant detection network trained on synthetic manuscript
          degradations (ink bleed, warping, damage).
    \item A stemma reconstruction algorithm using neural embeddings
          of textual variants, validated on both synthetic and real corpora.
\end{enumerate}

% ============================================================
% 2. RELATED WORK
% ============================================================
\section{Related Work}
\label{sec:related}

\subsection{Neural Rendering}
\label{sec:rw-rendering}

NeRF \parencite{mildenhall2021nerf} introduced coordinate-based MLPs for
novel view synthesis. Extensions include instant-NGP \parencite{muller2022instant},
3D Gaussian Splatting \parencite{kerbl20233d}, and differentiable
rasterization \parencite{laine2020diffrast}.

\subsection{Digital Humanities \& Critical Editions}
\label{sec:rw-dh}

Computational collation tools include CollateX \parencite{dekker2015collatex},
Juxta \parencite{juxta}, and the Versioning Machine \parencite{vmachine}.
Stemmatology uses cladistic methods \parencite{spencer2004cladistics}
and Bayesian approaches \parencite{barbrook2005bayesian}.

\subsection{Cross-Domain Applications}
\label{sec:rw-cross}

Recent work applies graphics to DH: document enhancement
\parencite{stamatopoulos2020document}, layout analysis
\parencite{pavlopoulos2023layout}, and HTR \parencite{bluche2017htr}.
Our work is the first to use differentiable rendering for collation.

% ============================================================
% 3. METHOD
% ============================================================
\section{Method}
\label{sec:method}

\subsection{Problem Formulation}
\label{sec:formulation}

Given manuscript images $\{I_k\}_{k=1}^N$ of the same text passage,
we seek the archetypal text $T^*$ and stemma $\mathcal{S}$.
We model each witness as a degraded rendering:

\begin{equation}\label{eq:degradation}
I_k = \mathcal{D}_k \circ \mathcal{R}(T^*) + \epsilon_k
\end{equation}

where $\mathcal{R}$ renders text to image, $\mathcal{D}_k$ applies
witness-specific degradations, and $\epsilon_k$ is noise.

\subsection{Differentiable Text Renderer}
\label{sec:renderer}

We implement $\mathcal{R}$ as a differentiable pipeline:

\begin{enumerate}
    \item \textbf{Text Layout}: Characters $\to$ positioned glyphs
    \item \textbf{Rasterization}: Glyphs $\to$ soft coverage masks
    \item \textbf{Degradation}: Physics-based ink/paper simulation
\end{enumerate}

\textbf{Soft Rasterization.} For each pixel $p$, coverage of glyph $g$:
\begin{equation}
\alpha_g(p) = \sigma\left( \frac{d_g(p)}{\tau} \right)
\end{equation}
where $d_g(p)$ is signed distance, $\tau$ temperature, $\sigma$ sigmoid.

\textbf{Ink Simulation.} Blot model via convolution:
\begin{equation}
I_{ink} = \alpha * k_{blot} + \mathcal{N}(0, \sigma_{paper})
\end{equation}
with kernel $k_{blot}$ modeling feathering.

\subsection{Variant Detection Network}
\label{sec:variant-net}

Architecture: U-Net with attention, trained to predict
variant heatmaps from paired manuscript crops.

\begin{algorithm}[htbp]
\caption{Variant Detection Training}\label{alg:variant}
\Input{Paired crops $(I_a, I_b)$, ground truth variants $V$}
\Output{Detector $D_\theta$}
\For{epoch = 1 to $E$}{
    Sample batch $\{(I_a^i, I_b^i, V^i)\}$\;
    $\hat{V}^i \gets D_\theta(I_a^i, I_b^i)$\;
    $\mathcal{L} \gets \text{BCE}(\hat{V}^i, V^i) + \lambda \text{Dice}(\hat{V}^i, V^i)$\;
    $\theta \gets \theta - \eta \nabla_\theta \mathcal{L}$\;
}
\Return $D_\theta$
\end{algorithm}

\subsection{Stemma Reconstruction}
\label{sec:stemma-rec}

We embed each witness reading $r_v(w_i)$ into a latent space
using a transformer encoder. Pairwise distances $d_{ij}$ define
edge weights for maximum parsimony:

\begin{equation}\label{eq:parsimony}
\mathcal{S}^* = \argmin_{\mathcal{S}} \sum_{(i,j) \in \mathcal{S}} d_{ij}
\end{equation}

Solved via neighbor-joining with bootstrap support.

% ============================================================
% 4. EXPERIMENTS
% ============================================================
\section{Experiments}
\label{sec:experiments}

\subsection{Synthetic Manuscript Corpus}
\label{sec:synthetic}

We generate 10,000 synthetic witnesses from 100 base texts
using our differentiable renderer with randomized degradations:
\begin{itemize}
    \item Ink: bleed $\sim \mathcal{U}(0, 3px)$, feather $\sim \mathcal{U}(0, 2px)$
    \item Paper: texture $\sim$ Perlin noise, stains $\sim$ Poisson disk
    \item Damage: tears, holes, fading $\sim$ procedural
    \item Layout: warp $\sim$ TPS, skew $\sim \mathcal{U}(-5^\circ, 5^\circ)$
\end{itemize}

\begin{table}[htbp]
\caption{Synthetic corpus statistics}\label{tab:synthetic}
\begin{tabular}{lcc}
\toprule
Parameter & Range & Distribution \\
\midrule
Base texts & 100 & Latin, Greek, Medieval \\
Witnesses per text & 50--200 & Uniform \\
Variants per text & 10--50 & Poisson($\lambda=25$) \\
Degradation types & 12 & Categorical \\
\bottomrule
\end{tabular}
\end{table}

\subsection{Collation Accuracy}
\label{sec:collation-acc}

\begin{table}[htbp]
\caption{Variant detection F1 scores}\label{tab:collation}
\begin{tabular}{lcccc}
\toprule
Method & Precision & Recall & F1 & Speed (img/s) \\
\midrule
CollateX (baseline) & 0.72 & 0.68 & 0.70 & 120 \\
HTR + Alignment & 0.78 & 0.74 & 0.76 & 15 \\
Ours (no pretrain) & 0.85 & 0.81 & 0.83 & 45 \\
Ours (pretrained) & \textbf{0.92} & \textbf{0.89} & \textbf{0.90} & 45 \\
\bottomrule
\end{tabular}
\end{table}

\subsection{Stemma Reconstruction}
\label{sec:stemma-exp}

\begin{table}[htbp]
\caption{Stemma accuracy (Robinson-Foulds distance)}\label{tab:stemma}
\begin{tabular}{lccc}
\toprule
Method & RF Distance $\downarrow$ & Bootstrap Support & Time \\
\midrule
Cladistic (PAUP) & 0.34 & 0.62 & 120s \\
Bayesian (MrBayes) & 0.28 & 0.71 & 3600s \\
Neighbor-Joining & 0.31 & 0.65 & 2s \\
Ours (embedding + NJ) & \textbf{0.22} & \textbf{0.83} & 5s \\
\bottomrule
\end{tabular}
\end{table}

\subsection{Real Manuscript Case Study}
\label{sec:case-study}

We apply our pipeline to the \emph{Codex Vaticanus} / \emph{Sinaiticus}
Gospel parallels (Matthew 1--5), 472 variation units, 4 witnesses.

\begin{figure}[htbp]
\centering
\begin{subfigure}[b]{0.48\linewidth}
\centering
\fbox{\parbox{4cm}{\centering [Manuscript crop A]}}
\caption{Vaticanus}\label{subfig:vat}
\end{subfigure}
\hfill
\begin{subfigure}[b]{0.48\linewidth}
\centering
\fbox{\parbox{4cm}{\centering [Manuscript crop B]}}
\caption{Sinaiticus}\label{subfig:sin}
\end{subfigure}
\caption{Manuscript crops with detected variants highlighted}\label{fig:manuscripts}
\end{figure}

Results: 94\% variant recall, reconstructed stemma matches
scholarly consensus with 89\% bootstrap support.

% ============================================================
% 5. DISCUSSION
% ============================================================
\section{Discussion}
\label{sec:discussion}

\subsection{Limitations}
\label{sec:limitations}

\begin{itemize}
    \item Requires paired training data (synthetic mitigates this)
    \item Assumes known text segmentation (line detection needed)
    \item Latin/Greek scripts only; CJK/Indic future work
    \item Computational cost: 45 img/s on RTX 4090
\end{itemize}

\subsection{Applications Beyond Manuscripts}
\label{sec:applications}

\begin{itemize}
    \item \textbf{Inscriptions}: 3D scanning + differentiable rendering for eroded stone
    \item \textbf{Palimpsests}: Multi-spectral layers as volumetric channels
    \item \textbf{Printed books}: Edition collation for early modern texts
    \item \textbf{Born-digital}: Version control for collaborative writing
\end{itemize}

\subsection{Ethical Considerations}
\label{sec:ethics}

Automated collation must not replace expert judgment.
Our tool outputs \emph{probabilistic} variants requiring
scholarly verification. We release model weights with
``human-in-the-loop'' licensing.

% ============================================================
% 6. CONCLUSION
% ============================================================
\section{Conclusion}
\label{sec:conclusion}

We unified neural rendering and critical edition through
differentiable text rendering. Our framework achieves
state-of-the-art collation and stemma reconstruction,
demonstrating that graphics techniques transfer productively
to Digital Humanities. Code and data at \url{github.com/example/neural-collation}.

% ============================================================
% ACKNOWLEDGMENTS
% ============================================================
\section*{Acknowledgments}

We thank the Berlin State Library for manuscript access.
Funded by DFG Grant SM 123/4-1 and JSPS KAKENHI 23K12345.

% ============================================================
% BIBLIOGRAPHY
% ============================================================
\ifdhformat
\printbibliography
\else
\bibliography{references}
\fi

% ============================================================
% APPENDICES (DH format only)
% ============================================================
\ifdhformat
\appendix
\section{Supplemental: Full Variant Tables}
\label{app:variants}

% Full apparatus would go here
See supplementary material for complete apparatus criticus.

\section{Supplemental: Stemma Visualization}
\label{app:stemma-full}

\begin{figure}[htbp]
\centering
\begin{tikzpicture}[
    node distance=1cm and 1.5cm,
    ms/.style={ellipse, draw, minimum width=1.5cm, minimum height=0.7cm, fill=blue!10, font=\itshape\small},
    lost/.style={ellipse, draw, dashed, minimum width=1.5cm, minimum height=0.7cm, fill=red!10, font=\itshape\small},
    arrow/.style={-Stealth, thick}
]
\node[lost] (O) {O};
\node[lost, below left=of O] (a) {$\alpha$};
\node[lost, below right=of O] (b) {$\beta$};
\node[ms, below left=of a] (A) {A};
\node[ms, below right=of a] (B) {B};
\node[ms, below left=of b] (C) {C};
\node[ms, below right=of b] (D) {D};
\draw[arrow] (O) -- (a); \draw[arrow] (O) -- (b);
\draw[arrow] (a) -- (A); \draw[arrow] (a) -- (B);
\draw[arrow] (b) -- (C); \draw[arrow] (b) -- (D);
\end{tikzpicture}
\caption{Full stemma with bootstrap values}\label{fig:stemma-full}
\end{figure}

\section{Supplemental: Degradation Parameters}
\label{app:degradation}

\begin{table}[htbp]
\caption{Full degradation parameter ranges}\label{tab:degradation-full}
\begin{tabular}{lll}
\toprule
Category & Parameter & Range \\
\midrule
Ink & bleed\_radius & $\mathcal{U}(0, 3)$ px \\
Ink & feather\_sigma & $\mathcal{U}(0, 2)$ px \\
Ink & intensity\_drop & $\mathcal{U}(0.7, 1.0)$ \\
Paper & perlin\_scale & $\mathcal{U}(0.01, 0.05)$ \\
Paper & stain\_count & $\text{Poisson}(3)$ \\
Paper & stain\_opacity & $\mathcal{U}(0.1, 0.4)$ \\
Damage & tear\_prob & 0.1 \\
Damage & hole\_prob & 0.05 \\
Damage & fade\_factor & $\mathcal{U}(0.5, 1.0)$ \\
Layout & warp\_magnitude & $\mathcal{U}(0, 10)$ px \\
Layout & skew\_angle & $\mathcal{U}(-5^\circ, 5^\circ)$ \\
\bottomrule
\end{tabular}
\end{table}

\fi

\end{document}